% =============================================================================
%  Comands.tex — Comandos y paquetes personalizados
%  Tesis: Transportes Anillares y su Importancia en la CDMX
%  Hernán Galileo Cabrera Garibaldi, UNAM 2026
%
%  NOTA: Este archivo se carga desde tesis.tex mediante % =============================================================================
%  Comands.tex — Comandos y paquetes personalizados
%  Tesis: Transportes Anillares y su Importancia en la CDMX
%  Hernán Galileo Cabrera Garibaldi, UNAM 2026
%
%  NOTA: Este archivo se carga desde tesis.tex mediante % =============================================================================
%  Comands.tex — Comandos y paquetes personalizados
%  Tesis: Transportes Anillares y su Importancia en la CDMX
%  Hernán Galileo Cabrera Garibaldi, UNAM 2026
%
%  NOTA: Este archivo se carga desde tesis.tex mediante % =============================================================================
%  Comands.tex — Comandos y paquetes personalizados
%  Tesis: Transportes Anillares y su Importancia en la CDMX
%  Hernán Galileo Cabrera Garibaldi, UNAM 2026
%
%  NOTA: Este archivo se carga desde tesis.tex mediante \input{Latex/Comands}
%        DESPUÉS de haber declarado todos los paquetes base. No usar
%        \documentclass ni \begin{document} aquí.
% =============================================================================


% =============================================================================
%  1. PAQUETES ADICIONALES
% =============================================================================

% --- Figuras y subfiguras ---
\usepackage{caption}           % Control fino de pies de figura y tabla
\usepackage{subcaption}        % Subfiguras (a), (b), (c)...
\usepackage{float}             % Opción [H] para fijar figuras en su lugar
\usepackage{wrapfig}           % Figuras con texto envolvente (útil en mapas pequeños)

% --- Tablas ---
\usepackage{booktabs}          % Líneas profesionales: \toprule, \midrule, \bottomrule
\usepackage{tabularx}          % Tablas con columna X de ancho ajustable
\usepackage{multirow}          % Celdas que abarcan múltiples filas

% --- Código fuente (Go, Python, SQL, JSON) ---
\usepackage{listings}          % Bloques de código con resaltado de sintaxis
\usepackage{xcolor}            % Colores personalizados para el resaltado

% --- Diagramas y gráficos vectoriales ---
\usepackage{tikz}              % Diagramas vectoriales (grafos, esquemas de red)
\usetikzlibrary{
  arrows.meta,                 % Puntas de flecha modernas
  positioning,                 % Posicionamiento relativo de nodos
  shapes.geometric,            % Formas: círculo, rectángulo, diamante, etc.
  fit,                         % Nodo que envuelve a otros
  backgrounds,                 % Capas de fondo en diagramas
  calc                         % Cálculos de coordenadas
}
\usepackage{pgfplots}          % Gráficas de datos (barras, líneas, dispersión)
\pgfplotsset{compat=1.18}

% --- Animaciones (GIFs exportados como PDF animado) ---
\usepackage{animate}


% =============================================================================
%  2. CONFIGURACIÓN DE COLORES INSTITUCIONALES
% =============================================================================

% Paleta UNAM
\definecolor{unamAzul}    {RGB}{0,  56, 147}   % Azul UNAM oficial
\definecolor{unamOro}     {RGB}{244,180,  0}   % Oro/amarillo UNAM
\definecolor{unamGrisClaro}{RGB}{240,240,240}  % Fondo gris claro para tablas

% Paleta de transporte (para mapas y diagramas de red)
\definecolor{colorMetro}    {RGB}{ 30, 136, 229}  % Azul metro
\definecolor{colorMetrobus} {RGB}{229,  57,  53}  % Rojo Metrobús
\definecolor{colorCablebus} {RGB}{ 67, 160,  71}  % Verde Cablebús
\definecolor{colorPeriferico}{RGB}{255, 143,   0} % Naranja Periférico Sur
\definecolor{colorPropuesto} {RGB}{142,  36, 170} % Morado corredor propuesto

% Paleta para código fuente
\definecolor{codeGray}   {RGB}{245,245,245}
\definecolor{codeComment}{RGB}{117,117,117}
\definecolor{codeKeyword}{RGB}{  0, 98, 179}
\definecolor{codeString} {RGB}{183, 28,  28}
\definecolor{codeNumber} {RGB}{ 46,125,  50}


% =============================================================================
%  3. CONFIGURACIÓN DE BLOQUES DE CÓDIGO (listings)
% =============================================================================

% --- Estilo base compartido ---
\lstdefinestyle{estiloBase}{
  backgroundcolor = \color{codeGray},
  basicstyle      = \ttfamily\small,
  breaklines      = true,
  breakatwhitespace = true,
  frame           = single,
  framerule       = 0.4pt,
  rulecolor       = \color{gray!50},
  numbers         = left,
  numberstyle     = \tiny\color{gray},
  numbersep       = 8pt,
  showstringspaces= false,
  tabsize         = 4,
  captionpos      = b,
  keepspaces      = true,
  commentstyle    = \color{codeComment}\itshape,
  stringstyle     = \color{codeString},
  keywordstyle    = \color{codeKeyword}\bfseries,
}

% --- Go (Golang) — para Apimetro / backend ---
\lstdefinelanguage{Go}{
  morekeywords = {
    func, package, import, var, const, type, struct, interface,
    return, if, else, for, range, switch, case, default, defer,
    go, chan, select, map, make, new, nil, true, false,
    string, int, int64, float64, bool, byte, error
  },
  sensitive    = true,
  morecomment  = [l]{//},
  morecomment  = [s]{/*}{*/},
  morestring   = [b]",
  morestring   = [b]`,
}

\lstdefinestyle{estiloGo}{
  style    = estiloBase,
  language = Go,
}

% --- Python — para análisis de grafos / GeoPandas ---
\lstdefinestyle{estiloPython}{
  style    = estiloBase,
  language = Python,
}

% --- SQL / PostGIS ---
\lstdefinestyle{estiloSQL}{
  style    = estiloBase,
  language = SQL,
}

% --- GeoJSON / JSON ---
\lstdefinelanguage{JSON}{
  morestring   = [b]",
  morekeywords = {true, false, null},
  sensitive    = false,
}

\lstdefinestyle{estiloJSON}{
  style    = estiloBase,
  language = JSON,
}

% Etiqueta por defecto para todos los listings
\renewcommand{\lstlistingname}{Código}
\renewcommand{\lstlistlistingname}{Índice de Código}


% =============================================================================
%  4. CONFIGURACIÓN DE FIGURAS Y SUBFIGURAS
% =============================================================================

% Subfiguras con etiqueta (a), (b), (c)... en lugar de (a (por el Comands original)
\renewcommand*{\thesubfigure}{(\alph{subfigure})}

% Estilo de pies de figura: fuente pequeña, etiqueta en negrita
\captionsetup{
  font       = small,
  labelfont  = bf,
  labelsep   = period,
  justification = raggedright,
  singlelinecheck = false
}

% Comando abreviado para fuente de figura
% Uso: \fuentefigura{García-López (2019, p.~3).}
\newcommand{\fuentefigura}[1]{%
  \par\vspace{2pt}%
  {\small\textit{Fuente: #1}}%
}


% =============================================================================
%  5. COMANDOS MATEMÁTICOS — Modelo de Transporte y Teoría de Grafos
% =============================================================================

% --- Notación de grafos ---
\newcommand{\grafo}{\mathcal{G}}          % Grafo G
\newcommand{\nodos}{\mathcal{V}}          % Conjunto de vértices/nodos V
\newcommand{\aristas}{\mathcal{E}}        % Conjunto de aristas E
\newcommand{\nodo}[1]{v_{#1}}             % Nodo individual: \nodo{i}
\newcommand{\arista}[2]{e_{#1#2}}         % Arista: \arista{i}{j}
\newcommand{\peso}[2]{w_{#1#2}}           % Peso de arista: \peso{i}{j}
\newcommand{\costo}[2]{c_{#1#2}}          % Costo total: \costo{i}{j}

% --- Función de costo total (modelo ponderado) ---
% Uso: \costototal{i}{j} genera c_{ij} = d_{ij} + \alpha·a_{ij} + \beta·\delta_{ij}
\newcommand{\costototal}[2]{%
  c_{#1#2} = d_{#1#2} + \alpha \cdot a_{#1#2} + \beta \cdot \delta_{#1#2}%
}

% --- Notación de caminos mínimos ---
\newcommand{\dijkstra}{\textsc{Dijkstra}}
\newcommand{\floyd}{\textsc{Floyd-Warshall}}
\newcommand{\kruskal}{\textsc{Kruskal}}

% --- Modelo de 4 etapas ---
\newcommand{\generacion}{T_i^{\text{gen}}}     % Viajes generados en zona i
\newcommand{\atraccion}{T_j^{\text{atr}}}      % Viajes atraídos en zona j
\newcommand{\distribucion}{T_{ij}}             % Viajes distribuidos i→j
\newcommand{\particion}{p_m}                   % Partición modal modo m

% --- Índices de accesibilidad ---
\newcommand{\accesibilidad}{A_i}               % Índice de accesibilidad zona i
\newcommand{\impedancia}{f(c_{ij})}            % Función de impedancia

% --- Vectores y matrices (heredado del Comands original) ---
\newcommand{\beqhalf}{%
  \noindent\begin{minipage}[c]{.5\linewidth}\begin{equation}}
\newcommand{\eeqhalf}{\end{equation}\end{minipage}}
\newcommand{\eqhalf}[1]{\beqhalf #1 \eeqhalf}


% =============================================================================
%  6. ENTORNO DE SUBECUACIONES ETIQUETADAS (heredado del Comands original)
% =============================================================================

\makeatletter
\newenvironment{taggedsubequations}[1]
 {%
   \addtocounter{equation}{-1}%
   \begin{subequations}%
   \def\@currentlabel{#1}%
   \renewcommand{\theequation}{#1\alph{equation}}%
 }
 {\end{subequations}}
\makeatother


% =============================================================================
%  7. COMANDOS DE UTILIDAD GENERAL
% =============================================================================

% Abreviaturas frecuentes en el texto
\newcommand{\cdmx}{Ciudad de México}
\newcommand{\zmvm}{Zona Metropolitana del Valle de México}
\newcommand{\brt}{\textsc{brt}}             % Bus Rapid Transit
\newcommand{\gis}{\textsc{gis}}             % Geographic Information System
\newcommand{\gtfs}{\textsc{gtfs}}           % General Transit Feed Specification
\newcommand{\api}{\textsc{api}}             % Application Programming Interface
\newcommand{\geojson}{GeoJSON}

% Nota al margen para revisión (solo en borrador — comentar antes de entregar)
% Uso: \pendiente{Completar con datos de campo}
\newcommand{\pendiente}[1]{%
  \marginpar{\footnotesize\color{red}\textbf{[Pendiente:} #1\textbf{]}}%
}

% Resaltado de término definido por primera vez
% Uso: \termino{fragmentación urbana}
\newcommand{\termino}[1]{\textit{#1}}
%        DESPUÉS de haber declarado todos los paquetes base. No usar
%        \documentclass ni \begin{document} aquí.
% =============================================================================


% =============================================================================
%  1. PAQUETES ADICIONALES
% =============================================================================

% --- Figuras y subfiguras ---
\usepackage{caption}           % Control fino de pies de figura y tabla
\usepackage{subcaption}        % Subfiguras (a), (b), (c)...
\usepackage{float}             % Opción [H] para fijar figuras en su lugar
\usepackage{wrapfig}           % Figuras con texto envolvente (útil en mapas pequeños)

% --- Tablas ---
\usepackage{booktabs}          % Líneas profesionales: \toprule, \midrule, \bottomrule
\usepackage{tabularx}          % Tablas con columna X de ancho ajustable
\usepackage{multirow}          % Celdas que abarcan múltiples filas

% --- Código fuente (Go, Python, SQL, JSON) ---
\usepackage{listings}          % Bloques de código con resaltado de sintaxis
\usepackage{xcolor}            % Colores personalizados para el resaltado

% --- Diagramas y gráficos vectoriales ---
\usepackage{tikz}              % Diagramas vectoriales (grafos, esquemas de red)
\usetikzlibrary{
  arrows.meta,                 % Puntas de flecha modernas
  positioning,                 % Posicionamiento relativo de nodos
  shapes.geometric,            % Formas: círculo, rectángulo, diamante, etc.
  fit,                         % Nodo que envuelve a otros
  backgrounds,                 % Capas de fondo en diagramas
  calc                         % Cálculos de coordenadas
}
\usepackage{pgfplots}          % Gráficas de datos (barras, líneas, dispersión)
\pgfplotsset{compat=1.18}

% --- Animaciones (GIFs exportados como PDF animado) ---
\usepackage{animate}


% =============================================================================
%  2. CONFIGURACIÓN DE COLORES INSTITUCIONALES
% =============================================================================

% Paleta UNAM
\definecolor{unamAzul}    {RGB}{0,  56, 147}   % Azul UNAM oficial
\definecolor{unamOro}     {RGB}{244,180,  0}   % Oro/amarillo UNAM
\definecolor{unamGrisClaro}{RGB}{240,240,240}  % Fondo gris claro para tablas

% Paleta de transporte (para mapas y diagramas de red)
\definecolor{colorMetro}    {RGB}{ 30, 136, 229}  % Azul metro
\definecolor{colorMetrobus} {RGB}{229,  57,  53}  % Rojo Metrobús
\definecolor{colorCablebus} {RGB}{ 67, 160,  71}  % Verde Cablebús
\definecolor{colorPeriferico}{RGB}{255, 143,   0} % Naranja Periférico Sur
\definecolor{colorPropuesto} {RGB}{142,  36, 170} % Morado corredor propuesto

% Paleta para código fuente
\definecolor{codeGray}   {RGB}{245,245,245}
\definecolor{codeComment}{RGB}{117,117,117}
\definecolor{codeKeyword}{RGB}{  0, 98, 179}
\definecolor{codeString} {RGB}{183, 28,  28}
\definecolor{codeNumber} {RGB}{ 46,125,  50}


% =============================================================================
%  3. CONFIGURACIÓN DE BLOQUES DE CÓDIGO (listings)
% =============================================================================

% --- Estilo base compartido ---
\lstdefinestyle{estiloBase}{
  backgroundcolor = \color{codeGray},
  basicstyle      = \ttfamily\small,
  breaklines      = true,
  breakatwhitespace = true,
  frame           = single,
  framerule       = 0.4pt,
  rulecolor       = \color{gray!50},
  numbers         = left,
  numberstyle     = \tiny\color{gray},
  numbersep       = 8pt,
  showstringspaces= false,
  tabsize         = 4,
  captionpos      = b,
  keepspaces      = true,
  commentstyle    = \color{codeComment}\itshape,
  stringstyle     = \color{codeString},
  keywordstyle    = \color{codeKeyword}\bfseries,
}

% --- Go (Golang) — para Apimetro / backend ---
\lstdefinelanguage{Go}{
  morekeywords = {
    func, package, import, var, const, type, struct, interface,
    return, if, else, for, range, switch, case, default, defer,
    go, chan, select, map, make, new, nil, true, false,
    string, int, int64, float64, bool, byte, error
  },
  sensitive    = true,
  morecomment  = [l]{//},
  morecomment  = [s]{/*}{*/},
  morestring   = [b]",
  morestring   = [b]`,
}

\lstdefinestyle{estiloGo}{
  style    = estiloBase,
  language = Go,
}

% --- Python — para análisis de grafos / GeoPandas ---
\lstdefinestyle{estiloPython}{
  style    = estiloBase,
  language = Python,
}

% --- SQL / PostGIS ---
\lstdefinestyle{estiloSQL}{
  style    = estiloBase,
  language = SQL,
}

% --- GeoJSON / JSON ---
\lstdefinelanguage{JSON}{
  morestring   = [b]",
  morekeywords = {true, false, null},
  sensitive    = false,
}

\lstdefinestyle{estiloJSON}{
  style    = estiloBase,
  language = JSON,
}

% Etiqueta por defecto para todos los listings
\renewcommand{\lstlistingname}{Código}
\renewcommand{\lstlistlistingname}{Índice de Código}


% =============================================================================
%  4. CONFIGURACIÓN DE FIGURAS Y SUBFIGURAS
% =============================================================================

% Subfiguras con etiqueta (a), (b), (c)... en lugar de (a (por el Comands original)
\renewcommand*{\thesubfigure}{(\alph{subfigure})}

% Estilo de pies de figura: fuente pequeña, etiqueta en negrita
\captionsetup{
  font       = small,
  labelfont  = bf,
  labelsep   = period,
  justification = raggedright,
  singlelinecheck = false
}

% Comando abreviado para fuente de figura
% Uso: \fuentefigura{García-López (2019, p.~3).}
\newcommand{\fuentefigura}[1]{%
  \par\vspace{2pt}%
  {\small\textit{Fuente: #1}}%
}


% =============================================================================
%  5. COMANDOS MATEMÁTICOS — Modelo de Transporte y Teoría de Grafos
% =============================================================================

% --- Notación de grafos ---
\newcommand{\grafo}{\mathcal{G}}          % Grafo G
\newcommand{\nodos}{\mathcal{V}}          % Conjunto de vértices/nodos V
\newcommand{\aristas}{\mathcal{E}}        % Conjunto de aristas E
\newcommand{\nodo}[1]{v_{#1}}             % Nodo individual: \nodo{i}
\newcommand{\arista}[2]{e_{#1#2}}         % Arista: \arista{i}{j}
\newcommand{\peso}[2]{w_{#1#2}}           % Peso de arista: \peso{i}{j}
\newcommand{\costo}[2]{c_{#1#2}}          % Costo total: \costo{i}{j}

% --- Función de costo total (modelo ponderado) ---
% Uso: \costototal{i}{j} genera c_{ij} = d_{ij} + \alpha·a_{ij} + \beta·\delta_{ij}
\newcommand{\costototal}[2]{%
  c_{#1#2} = d_{#1#2} + \alpha \cdot a_{#1#2} + \beta \cdot \delta_{#1#2}%
}

% --- Notación de caminos mínimos ---
\newcommand{\dijkstra}{\textsc{Dijkstra}}
\newcommand{\floyd}{\textsc{Floyd-Warshall}}
\newcommand{\kruskal}{\textsc{Kruskal}}

% --- Modelo de 4 etapas ---
\newcommand{\generacion}{T_i^{\text{gen}}}     % Viajes generados en zona i
\newcommand{\atraccion}{T_j^{\text{atr}}}      % Viajes atraídos en zona j
\newcommand{\distribucion}{T_{ij}}             % Viajes distribuidos i→j
\newcommand{\particion}{p_m}                   % Partición modal modo m

% --- Índices de accesibilidad ---
\newcommand{\accesibilidad}{A_i}               % Índice de accesibilidad zona i
\newcommand{\impedancia}{f(c_{ij})}            % Función de impedancia

% --- Vectores y matrices (heredado del Comands original) ---
\newcommand{\beqhalf}{%
  \noindent\begin{minipage}[c]{.5\linewidth}\begin{equation}}
\newcommand{\eeqhalf}{\end{equation}\end{minipage}}
\newcommand{\eqhalf}[1]{\beqhalf #1 \eeqhalf}


% =============================================================================
%  6. ENTORNO DE SUBECUACIONES ETIQUETADAS (heredado del Comands original)
% =============================================================================

\makeatletter
\newenvironment{taggedsubequations}[1]
 {%
   \addtocounter{equation}{-1}%
   \begin{subequations}%
   \def\@currentlabel{#1}%
   \renewcommand{\theequation}{#1\alph{equation}}%
 }
 {\end{subequations}}
\makeatother


% =============================================================================
%  7. COMANDOS DE UTILIDAD GENERAL
% =============================================================================

% Abreviaturas frecuentes en el texto
\newcommand{\cdmx}{Ciudad de México}
\newcommand{\zmvm}{Zona Metropolitana del Valle de México}
\newcommand{\brt}{\textsc{brt}}             % Bus Rapid Transit
\newcommand{\gis}{\textsc{gis}}             % Geographic Information System
\newcommand{\gtfs}{\textsc{gtfs}}           % General Transit Feed Specification
\newcommand{\api}{\textsc{api}}             % Application Programming Interface
\newcommand{\geojson}{GeoJSON}

% Nota al margen para revisión (solo en borrador — comentar antes de entregar)
% Uso: \pendiente{Completar con datos de campo}
\newcommand{\pendiente}[1]{%
  \marginpar{\footnotesize\color{red}\textbf{[Pendiente:} #1\textbf{]}}%
}

% Resaltado de término definido por primera vez
% Uso: \termino{fragmentación urbana}
\newcommand{\termino}[1]{\textit{#1}}
%        DESPUÉS de haber declarado todos los paquetes base. No usar
%        \documentclass ni \begin{document} aquí.
% =============================================================================


% =============================================================================
%  1. PAQUETES ADICIONALES
% =============================================================================

% --- Figuras y subfiguras ---
\usepackage{caption}           % Control fino de pies de figura y tabla
\usepackage{subcaption}        % Subfiguras (a), (b), (c)...
\usepackage{float}             % Opción [H] para fijar figuras en su lugar
\usepackage{wrapfig}           % Figuras con texto envolvente (útil en mapas pequeños)

% --- Tablas ---
\usepackage{booktabs}          % Líneas profesionales: \toprule, \midrule, \bottomrule
\usepackage{tabularx}          % Tablas con columna X de ancho ajustable
\usepackage{multirow}          % Celdas que abarcan múltiples filas

% --- Código fuente (Go, Python, SQL, JSON) ---
\usepackage{listings}          % Bloques de código con resaltado de sintaxis
\usepackage{xcolor}            % Colores personalizados para el resaltado

% --- Diagramas y gráficos vectoriales ---
\usepackage{tikz}              % Diagramas vectoriales (grafos, esquemas de red)
\usetikzlibrary{
  arrows.meta,                 % Puntas de flecha modernas
  positioning,                 % Posicionamiento relativo de nodos
  shapes.geometric,            % Formas: círculo, rectángulo, diamante, etc.
  fit,                         % Nodo que envuelve a otros
  backgrounds,                 % Capas de fondo en diagramas
  calc                         % Cálculos de coordenadas
}
\usepackage{pgfplots}          % Gráficas de datos (barras, líneas, dispersión)
\pgfplotsset{compat=1.18}

% --- Animaciones (GIFs exportados como PDF animado) ---
\usepackage{animate}


% =============================================================================
%  2. CONFIGURACIÓN DE COLORES INSTITUCIONALES
% =============================================================================

% Paleta UNAM
\definecolor{unamAzul}    {RGB}{0,  56, 147}   % Azul UNAM oficial
\definecolor{unamOro}     {RGB}{244,180,  0}   % Oro/amarillo UNAM
\definecolor{unamGrisClaro}{RGB}{240,240,240}  % Fondo gris claro para tablas

% Paleta de transporte (para mapas y diagramas de red)
\definecolor{colorMetro}    {RGB}{ 30, 136, 229}  % Azul metro
\definecolor{colorMetrobus} {RGB}{229,  57,  53}  % Rojo Metrobús
\definecolor{colorCablebus} {RGB}{ 67, 160,  71}  % Verde Cablebús
\definecolor{colorPeriferico}{RGB}{255, 143,   0} % Naranja Periférico Sur
\definecolor{colorPropuesto} {RGB}{142,  36, 170} % Morado corredor propuesto

% Paleta para código fuente
\definecolor{codeGray}   {RGB}{245,245,245}
\definecolor{codeComment}{RGB}{117,117,117}
\definecolor{codeKeyword}{RGB}{  0, 98, 179}
\definecolor{codeString} {RGB}{183, 28,  28}
\definecolor{codeNumber} {RGB}{ 46,125,  50}


% =============================================================================
%  3. CONFIGURACIÓN DE BLOQUES DE CÓDIGO (listings)
% =============================================================================

% --- Estilo base compartido ---
\lstdefinestyle{estiloBase}{
  backgroundcolor = \color{codeGray},
  basicstyle      = \ttfamily\small,
  breaklines      = true,
  breakatwhitespace = true,
  frame           = single,
  framerule       = 0.4pt,
  rulecolor       = \color{gray!50},
  numbers         = left,
  numberstyle     = \tiny\color{gray},
  numbersep       = 8pt,
  showstringspaces= false,
  tabsize         = 4,
  captionpos      = b,
  keepspaces      = true,
  commentstyle    = \color{codeComment}\itshape,
  stringstyle     = \color{codeString},
  keywordstyle    = \color{codeKeyword}\bfseries,
}

% --- Go (Golang) — para Apimetro / backend ---
\lstdefinelanguage{Go}{
  morekeywords = {
    func, package, import, var, const, type, struct, interface,
    return, if, else, for, range, switch, case, default, defer,
    go, chan, select, map, make, new, nil, true, false,
    string, int, int64, float64, bool, byte, error
  },
  sensitive    = true,
  morecomment  = [l]{//},
  morecomment  = [s]{/*}{*/},
  morestring   = [b]",
  morestring   = [b]`,
}

\lstdefinestyle{estiloGo}{
  style    = estiloBase,
  language = Go,
}

% --- Python — para análisis de grafos / GeoPandas ---
\lstdefinestyle{estiloPython}{
  style    = estiloBase,
  language = Python,
}

% --- SQL / PostGIS ---
\lstdefinestyle{estiloSQL}{
  style    = estiloBase,
  language = SQL,
}

% --- GeoJSON / JSON ---
\lstdefinelanguage{JSON}{
  morestring   = [b]",
  morekeywords = {true, false, null},
  sensitive    = false,
}

\lstdefinestyle{estiloJSON}{
  style    = estiloBase,
  language = JSON,
}

% Etiqueta por defecto para todos los listings
\renewcommand{\lstlistingname}{Código}
\renewcommand{\lstlistlistingname}{Índice de Código}


% =============================================================================
%  4. CONFIGURACIÓN DE FIGURAS Y SUBFIGURAS
% =============================================================================

% Subfiguras con etiqueta (a), (b), (c)... en lugar de (a (por el Comands original)
\renewcommand*{\thesubfigure}{(\alph{subfigure})}

% Estilo de pies de figura: fuente pequeña, etiqueta en negrita
\captionsetup{
  font       = small,
  labelfont  = bf,
  labelsep   = period,
  justification = raggedright,
  singlelinecheck = false
}

% Comando abreviado para fuente de figura
% Uso: \fuentefigura{García-López (2019, p.~3).}
\newcommand{\fuentefigura}[1]{%
  \par\vspace{2pt}%
  {\small\textit{Fuente: #1}}%
}


% =============================================================================
%  5. COMANDOS MATEMÁTICOS — Modelo de Transporte y Teoría de Grafos
% =============================================================================

% --- Notación de grafos ---
\newcommand{\grafo}{\mathcal{G}}          % Grafo G
\newcommand{\nodos}{\mathcal{V}}          % Conjunto de vértices/nodos V
\newcommand{\aristas}{\mathcal{E}}        % Conjunto de aristas E
\newcommand{\nodo}[1]{v_{#1}}             % Nodo individual: \nodo{i}
\newcommand{\arista}[2]{e_{#1#2}}         % Arista: \arista{i}{j}
\newcommand{\peso}[2]{w_{#1#2}}           % Peso de arista: \peso{i}{j}
\newcommand{\costo}[2]{c_{#1#2}}          % Costo total: \costo{i}{j}

% --- Función de costo total (modelo ponderado) ---
% Uso: \costototal{i}{j} genera c_{ij} = d_{ij} + \alpha·a_{ij} + \beta·\delta_{ij}
\newcommand{\costototal}[2]{%
  c_{#1#2} = d_{#1#2} + \alpha \cdot a_{#1#2} + \beta \cdot \delta_{#1#2}%
}

% --- Notación de caminos mínimos ---
\newcommand{\dijkstra}{\textsc{Dijkstra}}
\newcommand{\floyd}{\textsc{Floyd-Warshall}}
\newcommand{\kruskal}{\textsc{Kruskal}}

% --- Modelo de 4 etapas ---
\newcommand{\generacion}{T_i^{\text{gen}}}     % Viajes generados en zona i
\newcommand{\atraccion}{T_j^{\text{atr}}}      % Viajes atraídos en zona j
\newcommand{\distribucion}{T_{ij}}             % Viajes distribuidos i→j
\newcommand{\particion}{p_m}                   % Partición modal modo m

% --- Índices de accesibilidad ---
\newcommand{\accesibilidad}{A_i}               % Índice de accesibilidad zona i
\newcommand{\impedancia}{f(c_{ij})}            % Función de impedancia

% --- Vectores y matrices (heredado del Comands original) ---
\newcommand{\beqhalf}{%
  \noindent\begin{minipage}[c]{.5\linewidth}\begin{equation}}
\newcommand{\eeqhalf}{\end{equation}\end{minipage}}
\newcommand{\eqhalf}[1]{\beqhalf #1 \eeqhalf}


% =============================================================================
%  6. ENTORNO DE SUBECUACIONES ETIQUETADAS (heredado del Comands original)
% =============================================================================

\makeatletter
\newenvironment{taggedsubequations}[1]
 {%
   \addtocounter{equation}{-1}%
   \begin{subequations}%
   \def\@currentlabel{#1}%
   \renewcommand{\theequation}{#1\alph{equation}}%
 }
 {\end{subequations}}
\makeatother


% =============================================================================
%  7. COMANDOS DE UTILIDAD GENERAL
% =============================================================================

% Abreviaturas frecuentes en el texto
\newcommand{\cdmx}{Ciudad de México}
\newcommand{\zmvm}{Zona Metropolitana del Valle de México}
\newcommand{\brt}{\textsc{brt}}             % Bus Rapid Transit
\newcommand{\gis}{\textsc{gis}}             % Geographic Information System
\newcommand{\gtfs}{\textsc{gtfs}}           % General Transit Feed Specification
\newcommand{\api}{\textsc{api}}             % Application Programming Interface
\newcommand{\geojson}{GeoJSON}

% Nota al margen para revisión (solo en borrador — comentar antes de entregar)
% Uso: \pendiente{Completar con datos de campo}
\newcommand{\pendiente}[1]{%
  \marginpar{\footnotesize\color{red}\textbf{[Pendiente:} #1\textbf{]}}%
}

% Resaltado de término definido por primera vez
% Uso: \termino{fragmentación urbana}
\newcommand{\termino}[1]{\textit{#1}}
%        DESPUÉS de haber declarado todos los paquetes base. No usar
%        \documentclass ni \begin{document} aquí.
% =============================================================================


% =============================================================================
%  1. PAQUETES ADICIONALES
% =============================================================================

% --- Figuras y subfiguras ---
\usepackage{caption}           % Control fino de pies de figura y tabla
\usepackage{subcaption}        % Subfiguras (a), (b), (c)...
\usepackage{float}             % Opción [H] para fijar figuras en su lugar
\usepackage{wrapfig}           % Figuras con texto envolvente (útil en mapas pequeños)

% --- Tablas ---
\usepackage{booktabs}          % Líneas profesionales: \toprule, \midrule, \bottomrule
\usepackage{tabularx}          % Tablas con columna X de ancho ajustable
\usepackage{multirow}          % Celdas que abarcan múltiples filas

% --- Código fuente (Go, Python, SQL, JSON) ---
\usepackage{listings}          % Bloques de código con resaltado de sintaxis
\usepackage{xcolor}            % Colores personalizados para el resaltado

% --- Diagramas y gráficos vectoriales ---
\usepackage{tikz}              % Diagramas vectoriales (grafos, esquemas de red)
\usetikzlibrary{
  arrows.meta,                 % Puntas de flecha modernas
  positioning,                 % Posicionamiento relativo de nodos
  shapes.geometric,            % Formas: círculo, rectángulo, diamante, etc.
  fit,                         % Nodo que envuelve a otros
  backgrounds,                 % Capas de fondo en diagramas
  calc                         % Cálculos de coordenadas
}
\usepackage{pgfplots}          % Gráficas de datos (barras, líneas, dispersión)
\pgfplotsset{compat=1.18}

% --- Animaciones (GIFs exportados como PDF animado) ---
\usepackage{animate}


% =============================================================================
%  2. CONFIGURACIÓN DE COLORES INSTITUCIONALES
% =============================================================================

% Paleta UNAM
\definecolor{unamAzul}    {RGB}{0,  56, 147}   % Azul UNAM oficial
\definecolor{unamOro}     {RGB}{244,180,  0}   % Oro/amarillo UNAM
\definecolor{unamGrisClaro}{RGB}{240,240,240}  % Fondo gris claro para tablas

% Paleta de transporte (para mapas y diagramas de red)
\definecolor{colorMetro}    {RGB}{ 30, 136, 229}  % Azul metro
\definecolor{colorMetrobus} {RGB}{229,  57,  53}  % Rojo Metrobús
\definecolor{colorCablebus} {RGB}{ 67, 160,  71}  % Verde Cablebús
\definecolor{colorPeriferico}{RGB}{255, 143,   0} % Naranja Periférico Sur
\definecolor{colorPropuesto} {RGB}{142,  36, 170} % Morado corredor propuesto

% Paleta para código fuente
\definecolor{codeGray}   {RGB}{245,245,245}
\definecolor{codeComment}{RGB}{117,117,117}
\definecolor{codeKeyword}{RGB}{  0, 98, 179}
\definecolor{codeString} {RGB}{183, 28,  28}
\definecolor{codeNumber} {RGB}{ 46,125,  50}


% =============================================================================
%  3. CONFIGURACIÓN DE BLOQUES DE CÓDIGO (listings)
% =============================================================================

% --- Estilo base compartido ---
\lstdefinestyle{estiloBase}{
  backgroundcolor = \color{codeGray},
  basicstyle      = \ttfamily\small,
  breaklines      = true,
  breakatwhitespace = true,
  frame           = single,
  framerule       = 0.4pt,
  rulecolor       = \color{gray!50},
  numbers         = left,
  numberstyle     = \tiny\color{gray},
  numbersep       = 8pt,
  showstringspaces= false,
  tabsize         = 4,
  captionpos      = b,
  keepspaces      = true,
  commentstyle    = \color{codeComment}\itshape,
  stringstyle     = \color{codeString},
  keywordstyle    = \color{codeKeyword}\bfseries,
}

% --- Go (Golang) — para Apimetro / backend ---
\lstdefinelanguage{Go}{
  morekeywords = {
    func, package, import, var, const, type, struct, interface,
    return, if, else, for, range, switch, case, default, defer,
    go, chan, select, map, make, new, nil, true, false,
    string, int, int64, float64, bool, byte, error
  },
  sensitive    = true,
  morecomment  = [l]{//},
  morecomment  = [s]{/*}{*/},
  morestring   = [b]",
  morestring   = [b]`,
}

\lstdefinestyle{estiloGo}{
  style    = estiloBase,
  language = Go,
}

% --- Python — para análisis de grafos / GeoPandas ---
\lstdefinestyle{estiloPython}{
  style    = estiloBase,
  language = Python,
}

% --- SQL / PostGIS ---
\lstdefinestyle{estiloSQL}{
  style    = estiloBase,
  language = SQL,
}

% --- GeoJSON / JSON ---
\lstdefinelanguage{JSON}{
  morestring   = [b]",
  morekeywords = {true, false, null},
  sensitive    = false,
}

\lstdefinestyle{estiloJSON}{
  style    = estiloBase,
  language = JSON,
}

% Etiqueta por defecto para todos los listings
\renewcommand{\lstlistingname}{Código}
\renewcommand{\lstlistlistingname}{Índice de Código}


% =============================================================================
%  4. CONFIGURACIÓN DE FIGURAS Y SUBFIGURAS
% =============================================================================

% Subfiguras con etiqueta (a), (b), (c)... en lugar de (a (por el Comands original)
\renewcommand*{\thesubfigure}{(\alph{subfigure})}

% Estilo de pies de figura: fuente pequeña, etiqueta en negrita
\captionsetup{
  font       = small,
  labelfont  = bf,
  labelsep   = period,
  justification = raggedright,
  singlelinecheck = false
}

% Comando abreviado para fuente de figura
% Uso: \fuentefigura{García-López (2019, p.~3).}
\newcommand{\fuentefigura}[1]{%
  \par\vspace{2pt}%
  {\small\textit{Fuente: #1}}%
}


% =============================================================================
%  5. COMANDOS MATEMÁTICOS — Modelo de Transporte y Teoría de Grafos
% =============================================================================

% --- Notación de grafos ---
\newcommand{\grafo}{\mathcal{G}}          % Grafo G
\newcommand{\nodos}{\mathcal{V}}          % Conjunto de vértices/nodos V
\newcommand{\aristas}{\mathcal{E}}        % Conjunto de aristas E
\newcommand{\nodo}[1]{v_{#1}}             % Nodo individual: \nodo{i}
\newcommand{\arista}[2]{e_{#1#2}}         % Arista: \arista{i}{j}
\newcommand{\peso}[2]{w_{#1#2}}           % Peso de arista: \peso{i}{j}
\newcommand{\costo}[2]{c_{#1#2}}          % Costo total: \costo{i}{j}

% --- Función de costo total (modelo ponderado) ---
% Uso: \costototal{i}{j} genera c_{ij} = d_{ij} + \alpha·a_{ij} + \beta·\delta_{ij}
\newcommand{\costototal}[2]{%
  c_{#1#2} = d_{#1#2} + \alpha \cdot a_{#1#2} + \beta \cdot \delta_{#1#2}%
}

% --- Notación de caminos mínimos ---
\newcommand{\dijkstra}{\textsc{Dijkstra}}
\newcommand{\floyd}{\textsc{Floyd-Warshall}}
\newcommand{\kruskal}{\textsc{Kruskal}}

% --- Modelo de 4 etapas ---
\newcommand{\generacion}{T_i^{\text{gen}}}     % Viajes generados en zona i
\newcommand{\atraccion}{T_j^{\text{atr}}}      % Viajes atraídos en zona j
\newcommand{\distribucion}{T_{ij}}             % Viajes distribuidos i→j
\newcommand{\particion}{p_m}                   % Partición modal modo m

% --- Índices de accesibilidad ---
\newcommand{\accesibilidad}{A_i}               % Índice de accesibilidad zona i
\newcommand{\impedancia}{f(c_{ij})}            % Función de impedancia

% --- Vectores y matrices (heredado del Comands original) ---
\newcommand{\beqhalf}{%
  \noindent\begin{minipage}[c]{.5\linewidth}\begin{equation}}
\newcommand{\eeqhalf}{\end{equation}\end{minipage}}
\newcommand{\eqhalf}[1]{\beqhalf #1 \eeqhalf}


% =============================================================================
%  6. ENTORNO DE SUBECUACIONES ETIQUETADAS (heredado del Comands original)
% =============================================================================

\makeatletter
\newenvironment{taggedsubequations}[1]
 {%
   \addtocounter{equation}{-1}%
   \begin{subequations}%
   \def\@currentlabel{#1}%
   \renewcommand{\theequation}{#1\alph{equation}}%
 }
 {\end{subequations}}
\makeatother


% =============================================================================
%  7. COMANDOS DE UTILIDAD GENERAL
% =============================================================================

% Abreviaturas frecuentes en el texto
\newcommand{\cdmx}{Ciudad de México}
\newcommand{\zmvm}{Zona Metropolitana del Valle de México}
\newcommand{\brt}{\textsc{brt}}             % Bus Rapid Transit
\newcommand{\gis}{\textsc{gis}}             % Geographic Information System
\newcommand{\gtfs}{\textsc{gtfs}}           % General Transit Feed Specification
\newcommand{\api}{\textsc{api}}             % Application Programming Interface
\newcommand{\geojson}{GeoJSON}

% Nota al margen para revisión (solo en borrador — comentar antes de entregar)
% Uso: \pendiente{Completar con datos de campo}
\newcommand{\pendiente}[1]{%
  \marginpar{\footnotesize\color{red}\textbf{[Pendiente:} #1\textbf{]}}%
}

% Resaltado de término definido por primera vez
% Uso: \termino{fragmentación urbana}
\newcommand{\termino}[1]{\textit{#1}}